\section{Question 5: actes et conditions de la contrefaçon}

\subsection{Cas du droit national français}

Les actes de contrefaçon sont définis par L615-1 CPI renvoyant notamment à
L613-3 et L613-4 CPI.

\subsubsection{Dimitrios France}

\textbf{Element légal.} Dimitrios France pourrait avoir commis les actes suivants selon L613-3(a) CPI
au regard des revendications de produit:
\begin{itemize}
    \item fabrication du produit ("assemblage des pièces et intégration du
    disque de fixation" énoncé de la question 5, "dispositif White fabriqué
    à Strasbourg par la société Dimitrios France" énoncé page 1);
    \item offre du produit ("ses clients, services hospitaliers ou sociétés
    spécialisées dans les soins à domicile, sont sollicités par la multinationale
    Dimitrios qui leur offre le dispositif White" énoncé page 1);
    \item mise dans le commerce du produit;
    \item exportation du produit ("Dimitrios France a également des locaux
    au Havre depuis lesquels elle expédie le dispositif White vers ses filiales
    dont Dimitrios NL à Rotterdam." énoncé page 1);
    \item transbordement du produit (déduit de fait que le Havre est un port)
    \item détention en vue de l'offre, de la mise dans le commerce, 
    de l'exportation, du transbordement du produit.\\
\end{itemize}

\textbf{Element moral.} L615-1 al.3 CPI dispose que la fabrication du produit contrefaisant, son exportation,
son transbordement, sa détention en vue de l'exportation du transbordement n'admettent
aucune condition de responsabilité.\\

L615-1 al.3 CPI dispose que l'offre, la mise dans le commerce et la détention en vue
de l'offre et de la mise dans le commerce du produit admettent une condition de
responsabilité, à savoir la connaissance de cause. Toutefois, par exception
selon L615-1 al. 3, Dimitrios France, étant fabricant
du produit contrefaisant, ne bénéficie pas de condition de responsabilité. Même si
Dimitrios France réfute être fabricant du produit contrefaisant, il semble qu'un
autre argument s'oppose au bénéfice de la connaissance de cause. En effet, il paraît
que le domaine des consommables médicaux est restreint: "M. Mathis, un commercial
de la société Dimitrios France, a présenté le dispositif White comme une alternative
au dispositif Vesper. M. Mathis était auparavant commercial dans le domaine des
consommables médicaux.". Ainsi, en application du Pourvoi n°99-16.926, la
connaissance de cause est de surcroît présumée.\\

\textbf{Element légal.} Dimitrios France pourrait avoir commis les actes suivants selon L613-3(b) CPI
au regard de la revendication de procédé:
\begin{itemize}
    \item offre d'utilisation du procédé ("salon
    organisé à Paris, [...] un commercial de la société Dimitrios France, a 
    présenté le dispositif White comme une alternative au dispositif Vesper
    " énoncé page 2): offre en France d'utilisation en France.\\
\end{itemize}

\textbf{Element moral.} L613-3(b) CPI dispose que l'offre de l'utilisation du procédé objet du brevet admet
une condition de responsabilité, à savoir si le tiers sait ou les circonstances
rendent évident. De toute évidence, le tiers sait puisque M. Mathis a présenté
le dispositif White comme une alternative du dispositif Vesper.\\

\textbf{Conclusion: En cas de reproduction des revendications de la partie française du brevet européen,
la responsabilité civile de Dimitrios France est engagée.}

\subsubsection{Dimitrios NL}

Dimitrios NL ne peut pas être inculpée de contrefaçon devant le Tribunal Judiciaire
de Paris des sociétés.\\

\textbf{Conclusion: En cas de reproduction des revendications de la partie française du brevet européen,
la responsabilité civile de Dimitrios NL n'est pas engagée.}

\subsubsection{Skyfleed}

\textbf{Element légal.} Skyfleed pourrait avoir commis un acte de fourniture de moyen selon L613-4(1)CPI :
\begin{itemize}
    \item il y a livraison sur le territoire français ("Skyfleed livre en France
    à Dimitrios France" énoncé de la question 5);
    \item Skyfleed est une personne autre que celles habilitées à exploiter
    l'invention brevetée;
    \item il y a des moyens de mise en œuvre de l'invention se rapportant à un 
    élément essentiel de l'invention brevetée ("sous la forme
    de pièces détachées, la viole/écrou et la tige filetée du dispositif White"
    énoncé de la question 5) car l'invention ne pourrait être mise en œuvre
    sans ces pièces, les moyens faisant partie d'une combinaison (cf. Pourvoi
    n°15-29.378).\\
\end{itemize}

Toutefois, les moyens ne sont pas aptes ou destinés à mettre en œuvre l'invention.
En effet, l'objet litigieux permet d'injecter un fluide mais en aucun cas d'aspirer
un fluide.\\

\textbf{Element moral.} L613-4(1) CPI dispose que la contrefaçon par fourniture de moyens
admet une condition de responsabilité: le tiers sait ou les circonstances rendent
évident. Il n'y a pas de preuve que les circonstances rendent évident que ces moyens sont aptes et destinés 
à cette mise en œuvre.\\

\textbf{Conclusion: En cas de reproduction des revendications de la partie française du brevet européen,
la responsabilité civile de skyfleeed n'est pas engagée.}

\subsection{Cas de la JUB}

\subsubsection{Dimitrios France}

\textbf{Element légal.} Dimitrios France pourrait avoir commis les actes suivants selon A25(a) AJUb
au regard des revendications de produit:
\begin{itemize}
    \item fabriquer le produit qui fait l'objet du brevet;
    \item offrir le produit;
    \item mettre sur le marché le produit;
    \item détenir le produit pour l'offrir et le mettre sur le marché.\\
\end{itemize}

\textbf{Element moral.} Il n'y pas de condition de responsabilité pour les actes selon A25(a).\\

\textbf{Element légal.} Dimitrios France pourrait avoir commis les actes suivants selon A25(b) AJUB
au regard de la revendication de procédé:
\begin{itemize}
    \item offrir l'utilisation du procédé sur le territoire des Etats membres contractants dans lesquels
    le brevet produit ses effets.\\
\end{itemize}

\textbf{Element moral.} Même raisonnement que pour le cas du droit national français.\\

\textbf{Conclusion: Dimitrios France est susceptible de faire l'objet d'une condamnation pour contrefaçon
par la JUB des sociétés.}

\subsubsection{Dimitrios NL}

\textbf{Element légal.} Dimitrios NL pourrait avoir commis les actes suivants selon A25(a) AJUb
au regard des revendications de produit:
\begin{itemize}
    \item offrir le produit objet du brevet ("la filiale Dimitrios NL vend ensuite le dispositif White
    à différents services hospitaliers ou sociétés spécialisées situées aux Pays-Bas qui utilisent ce
    dispositif" énoncé page 1);
    \item mettre sur le marché le produit objet du brevet;
    \item utiliser le produit objet du brevet ("Dimitrios NL utilise également le dispositif White"
    énoncé page 1);
    \item détenir le produit pour l'offrir, le mettre sur le marché et l'utiliser; 
\end{itemize}

\textbf{Element moral.} Il n'y pas de condition de responsabilité pour les actes selon A25(a).\\

\textbf{Element légal.} Dimitrios France pourrait avoir commis les actes suivants selon A25(b) AJUB
au regard de la revendication de procédé:
\begin{itemize}
    \item utiliser le procédé qui fait l'objet du brevet;
    \item offrir l'utilisation du procédé sur le territoire des Etats membres contractants dans lesquels
    le brevet produit ses effets.\\
\end{itemize}

\textbf{Element moral.} Il n'y a pas de condition de responsabilité pour utiliser le procédé objet du brevet.
Pour l'offre d'utilisation du procédé, même raisonnement que pour le cas du droit national français.\\

\textbf{Conclusion: Skyfleed est susceptible de faire l'objet d'une condamnation pour contrefaçon
par la JUB des sociétés.}

\subsubsection{Skyfleed}

\textbf{Element légal.} Skyfleed pourrait avoir commis un acte de fourniture de moyen selon A26 AJUB.
Le même raisonnement que dans le cas du droit national français s'applique.\\

\textbf{Element moral.} Même raisonnement que pour le cas du droit national français.\\

\textbf{Conclusion: Skyfleed n'est pas susceptible de faire l'objet d'une condamnation pour contrefaçon
par la JUB des sociétés.}